% ============================================================================
% Appendix: Local Optimality of N_bonds = 3 in the Pinning Model
% ============================================================================
% Status: [Dc] within local pairwise-additive pinning model
% Scope: This is NOT a topological proof from S_EDC.
%        It is an optimality result within a declared effective energy model.
% Phase 1 / R4 deliverable (see PHASE1_PLAN_REVISED.md v3.1)
% ============================================================================

\chapter{Local Optimality of $N_{\text{bonds}} = 3$}
\label{app:nbonds}

\section*{Scope and Epistemic Status}

This appendix derives $N_{\text{bonds}} = 3$ as a \textbf{local optimality
result within the pairwise-additive pinning model}. It is not a direct
derivation from $S_{\text{EDC}}$, nor a general topological proof. The
result operates within the pinning Hamiltonian $H_{\text{pin}}$ defined in
Chapter~\ref{ch:sigma_to_K}, which itself contains postulated assumptions
about pairwise additivity and locality.

\textbf{Intended epistemic tag:} [Dc] within the local pairwise-additive
pinning model. A full [Der] promotion would require deriving $H_{\text{pin}}$
from $S_{\text{EDC}}$, which is beyond the scope of this appendix.

% ============================================================================

\begin{theorem}[Local Optimality of $N_{\text{bonds}} = 3$]
\label{thm:nbonds_optimality}
Within the pairwise-additive pinning model, for two Y-junctions sharing a
contact face in the $\Msix$ lattice, the energy-minimizing configuration
has exactly $N_{\text{bonds}} = 3$ contact pairs.
\end{theorem}

The proof proceeds via three lemmas.

% ============================================================================
\section{Lemma 1: Upper Bound}
% ============================================================================

\begin{lemma}[Upper Bound on Bond Count]
\label{lem:upper_bound}
Within the contact model, $N_{\text{bonds}} \leq 3$.
\end{lemma}

\begin{proof}
Each Y-junction has exactly 3 arms, by the $\Zthree$ Steiner geometry
established in Chapter~\ref{ch:anchor}~[Der].

Within the contact model, each arm can support at most one bond. This is
a geometric constraint within the model: a bond is defined as a pairwise
arm-to-arm contact between one arm of $J_1$ and one arm of $J_2$. Two
worldsheet segments sharing the same contact face are treated as a single
bond, not multiple independent bonds. This one-bond-per-arm constraint
follows from the model's definition of bond as a pairwise contact, but
has not been independently derived as a physical saturation theorem from
the full 5D action.

Since each junction contributes at most 3 arms, and each arm supports at
most one bond:
\begin{equation}
    N_{\text{bonds}} \leq \min(3, 3) = 3.
    \label{eq:nbonds_upper}
\end{equation}
\end{proof}

\textbf{Status:} [Der] for the arm count; contact-model constraint for
the one-bond-per-arm rule.

% ============================================================================
\section{Lemma 2: Energetic Preference for Saturation}
% ============================================================================

\begin{lemma}[Energetic Preference for Full Saturation]
\label{lem:saturation}
Within the local pairwise-additive pinning model, the total contact energy
$E(n)$ for $n$ bonded pairs ($0 \leq n \leq 3$) is minimized at $n = 3$.
\end{lemma}

\begin{proof}
Define the contact energy for $n$ bonded pairs:
\begin{equation}
    E(n) = n \times \Kpin + (3 - n) \times E_{\text{free}}
    \label{eq:contact_energy}
\end{equation}
where $\Kpin < 0$ (binding, from Chapter~\ref{ch:sigma_to_K}~[Dc]) and
$E_{\text{free}} = 0$ (free arm, reference level).

This energy model requires the following assumptions, all of which are
structurally necessary for the monotonicity argument:

\begin{enumerate}
    \item \textbf{Pairwise additivity}~[P]: The total energy is the sum of
          individual pair contributions. There are no irreducible three-body
          or higher-order terms.
    \item \textbf{Locality}~[P]: No long-range inter-arm coupling exists
          beyond the nearest contact. Arms that are not directly bonded do
          not interact.
    \item \textbf{No frustration penalty}~[P]: Adding a bond does not
          create geometric strain that offsets the binding energy gained.
          The junction geometry accommodates all three bonds without
          distortion cost.
    \item \textbf{No multi-arm nonlocal coupling}~[P]: There is no
          collective effect whereby the engagement of one arm pair modifies
          the binding energy available to another arm pair.
\end{enumerate}

Under these four assumptions, $E(n) = n \times \Kpin$ is monotonically
decreasing in $n$ (since $\Kpin < 0$), and is therefore minimized at
$n = 3$:
\begin{equation}
    E(3) = 3\Kpin < 2\Kpin = E(2) < \Kpin = E(1) < 0 = E(0).
    \label{eq:monotonic}
\end{equation}
\end{proof}

\begin{remark}[Monotonicity Failure Modes]
\label{rem:failure_modes}
If any of assumptions 1--4 fail, monotonicity in $n$ can fail. For example:
\begin{itemize}
    \item If frustration penalties grow faster than $|\Kpin|$ with increasing
          $n$, the minimum could shift to $n < 3$.
    \item If multi-arm coupling introduces a repulsive three-body term, the
          energy at $n = 3$ could exceed the energy at $n = 2$.
    \item If geometric strain from simultaneous three-arm engagement is
          large enough, $n = 2$ could be globally preferred.
\end{itemize}
None of these scenarios is excluded by the current argument. They would
require a different physical mechanism and can only be tested by explicit
computation from the 5D action.
\end{remark}

\textbf{Status:} [Dc] within the declared pairwise-additive pinning model.
This is not a general full-theory result.

% ============================================================================
\section{Lemma 3: Uniqueness of Pairing}
% ============================================================================

\begin{lemma}[Uniqueness up to $\Zthree$]
\label{lem:uniqueness}
Within the restricted geometric contact model (fixed Y-junction orientations,
fixed arm directions), the $n = 3$ bond pairing is unique up to a $\Zthree$
rotation.
\end{lemma}

\begin{proof}
The three arms of $J_1$ are directed at angles $\{0°, 120°, 240°\}$.
The three arms of $J_2$, facing $J_1$, are at complementary angles
$\{180°, 300°, 60°\}$.

A bond pairing is a bijection $\pi: \{a_1, a_2, a_3\} \to \{b_1, b_2, b_3\}$
assigning each arm of $J_1$ to an arm of $J_2$. The angular mismatch for
a pairing $\pi$ is:
\begin{equation}
    \Delta\theta(\pi) = \sum_{i=1}^{3}
        \bigl|\alpha_i^{(J_1)} - \alpha_{\pi(i)}^{(J_2)} - 180°\bigr|
    \label{eq:mismatch}
\end{equation}
where the subtraction of $180°$ accounts for the face-to-face orientation.

The identity pairing $\pi = \text{id}$ (arm$_i \leftrightarrow$ arm$_i'$)
has $\Delta\theta = 0$---perfect complementary alignment. Any non-identity
permutation $\pi \in S_3 \setminus \{\text{id}\}$ creates nonzero mismatch
for at least one pair. By the Steiner optimality principle
(Chapter~\ref{ch:anchor}~[Der]), nonzero angular mismatch increases the
brane area and hence the energy.

The minimum-energy pairing is therefore $\pi = \text{id}$, which is unique
within each $\Zthree$ orbit (the $\Zthree$ symmetry permutes the arm labels
cyclically, giving equivalent configurations).
\end{proof}

\textbf{Scope:} This uniqueness holds within the restricted geometric contact
model---fixed Y-junction orientations, fixed arm directions, and the mismatch
energy defined above. It is \emph{not} uniqueness in the full configuration
space of all admissible 5D brane deformations, where additional degrees of
freedom (arm bending, junction displacement, brane twisting) could modify
the energy landscape.

% ============================================================================
\section{Corollary: Deuterium Binding Energy}
% ============================================================================

\begin{corollary}[Binding Energy of the A=2 System]
\label{cor:Bd}
Within the local pairwise-additive pinning model,
\begin{equation}
    B_d = 3 \times \Kpin.
    \label{eq:Bd_derived}
\end{equation}
\end{corollary}

\begin{proof}
By Lemma~\ref{lem:upper_bound}, $N_{\text{bonds}} \leq 3$.
By Lemma~\ref{lem:saturation}, the energy minimum is at $n = 3$.
Therefore $N_{\text{bonds}} = 3$, and $B_d = 3\Kpin$ follows from the
linear pinning model (Eq.~\ref{eq:Bd_model} in Chapter~\ref{ch:deuterium}).
\end{proof}

\textbf{Numerical comparison:}
With $\Kpin \approx 0.74\MeV$~[Dc] (Chapter~\ref{ch:sigma_to_K}):
\begin{equation}
    B_d^{\text{(EDC)}} = 3 \times 0.74\MeV = 2.22\MeV.
\end{equation}
The baseline datum is $B_d^{\text{exp}} = 2.224\MeV$~[BL]. Agreement: $< 1\%$.

This numerical comparison is a consistency check, not derivation evidence.
The close agreement does not elevate the epistemic status of the model
assumptions.

% ============================================================================
\section{Residual Open Items}
% ============================================================================

The following items remain postulated~[P] and are not resolved by this
appendix:

\begin{enumerate}
    \item \textbf{Pairwise additivity of $H_{\text{pin}}$}~[P]: The pinning
          Hamiltonian is assumed to decompose as a sum of pair terms. Deriving
          this from $S_{\text{EDC}}$ requires completing the Put~C corridor
          (5D action $\to$ effective Hamiltonian reduction).

    \item \textbf{Locality (no long-range coupling)}~[P]: Arms not in direct
          contact are assumed non-interacting. In the full 5D theory, brane
          tension could mediate indirect coupling.

    \item \textbf{Frustration-free geometry}~[P]: The simultaneous engagement
          of all three arm pairs is assumed to incur no geometric strain cost.
          Verification requires explicit 5D embedding computation.

    \item \textbf{Absence of multi-arm collective effects}~[P]: No three-body
          or higher-order terms. This could be tested by computing the full
          contact energy from the brane action for the two-junction system.

    \item \textbf{Contact saturation (one bond per arm)}~[model constraint]:
          Used as a model definition in Lemma~\ref{lem:upper_bound}. A
          derivation from the 5D action could reveal whether multiple
          independent contacts per arm are physically possible.
\end{enumerate}

Upgrading any of these requires explicit computation from the 5D action,
which is the scope of the Put~C corridor work (Phase~1, Step~2/R3).